% !TEX root = ../main.tex
\chapter{一つの応答をたどる}
\label{chap:practice}

\begin{chapterabstract}
モデルが作った応答はHTTPの応答となり、ネットワークを戻り、端末のメモリへ届きます。ブラウザが画面を書き換えると、冒頭で送信ボタンを押した利用者に「明日の午後に到着する予定です」という文が見えます。止めていた数秒が、ようやく動き終えました。

この章では同じ旅を、今度は物理からアプリまで一息でたどります。そして、もし応答が遅い、届かない、メモリが増える、内容が誤っているなら、旅のどこへ戻るべきかを考えます。測定や記録によって確かめた情報を観測事実、現象を説明するため一時的に置く候補を仮説と呼びます。仮説が誤りなら現れないはずの事実を探して確かめることが反証です。分野名から原因を当てず、観測事実を境界へ置き、反証できる仮説へ変えるところで物語を閉じます。
\end{chapterabstract}

\begin{learninggoals}
  \item 一要求の観測点を、コード、実行、通信、モデル、描画へ配置できます。
  \item 症状から一つの原因を断定せず、反証可能な仮説へ分解できます。
  \item 遅延、メモリ、通信、描画、AI品質の調査手順を別々に組み立てられます。
\end{learninggoals}

\chapterimage{ch09-observability.png}{一要求の各境界に観測点を置き、異常箇所を切り分ける概念図}{fig:ch09-observability}


\section{リクエストの全体像}
\label{sec:07-practice-01}

利用者の画面へ応答が現れたところから、旅全体を一枚に戻します。上から見れば一つの要求と応答ですが、下へ降りるほど、意味を扱う層は計算、記憶、配送を担う層へ仕事を委ねています。どちら向きにもたどれる地図を先に持ちます。

\begingroup
  \centering
  \begin{tikzpicture}[node distance=2mm,every node/.append style={font=\tiny,minimum width=17mm,inner sep=1mm}]
    \node[accentbox] (user) {利用者の操作\\URL・API};
    \node[box,right=of user] (net) {通信\\DNS・HTTP};
    \node[box,right=of net] (app) {アプリ\\解析・業務処理};
    \node[box,right=of app] (runtime) {実行環境\\プロセス・OS};
    \node[box,right=of runtime] (machine) {機械\\CPU・メモリ};
    \node[accentbox,right=of machine] (signal) {物理的実体\\電圧・ビット};
    \draw[flow] (user)--(net); \draw[flow] (net)--(app); \draw[flow] (app)--(runtime);
    \draw[flow] (runtime)--(machine); \draw[flow] (machine)--(signal);
  \end{tikzpicture}
  \vspace{3mm}
\begingroup
\paragraph{上位層の役割}
        利用者の目的に近い意味を扱います。下位層の詳細を抽象化し、要求、データ、業務規則として表現します。

\par
\endgroup
\medskip
\begingroup
\paragraph{下位層の役割}
        上位層が使う計算、記憶、配送を実現します。障害時には抽象化の下へ一段ずつ降りて観測します。

\par
\endgroup
\par
\endgroup


\section{コードから実行まで}
\label{sec:07-practice-02}

地図の最下部から、もう一度上へ進みます。物理的な回路が命令を実行できても、その命令はソースコードの変換によって作られ、OSがプロセスとして起動したときに初めて要求を待てます。ソースコードと必要な部品を検査・変換して実行可能な成果物を作る段階が\term{ビルド時}です。OSがそれをプロセスにする段階が\term{起動時}です。プロセスが要求を処理する段階が\term{実行時}です。

\begingroup
  \centering
  \begin{tikzpicture}[node distance=4mm,every node/.append style={font=\scriptsize}]
    \node[accentbox] (src) {ソースコード};
    \node[box,right=of src] (front) {字句・構文解析\\AST};
    \node[box,right=of front] (code) {コード生成\\機械語};
    \node[box,right=of code] (exe) {実行可能形式};
    \draw[flow] (src)--(front); \draw[flow] (front)--(code); \draw[flow] (code)--(exe);
    \node[box,below=8mm of src] (loader) {OSのローダー};
    \node[box,right=of loader] (proc) {プロセス\\仮想アドレス空間};
    \node[box,right=of proc] (cpu) {CPUが命令を\\フェッチ・実行};
    \node[accentbox,right=of cpu] (sys) {システムコール\\入出力};
    \draw[flow] (exe.south) to[bend right=18] (loader.north);
    \draw[flow] (loader)--(proc); \draw[flow] (proc)--(cpu); \draw[flow] (cpu)--(sys);
  \end{tikzpicture}
  \vspace{5mm}
\begingroup
\paragraph{ビルド時}
        文法と型を検査し、必要なライブラリなどほかの部品との依存関係を解決して、対象の命令セットで実行できる形式を作ります。

\par
\endgroup
\medskip
\begingroup
\paragraph{起動時}
        OSがコードとデータを仮想メモリへ対応づけ、レジスタ状態、開いたファイル、権限など、実行を再開するための一式を作ります。

\par
\endgroup
\medskip
\begingroup
\paragraph{実行時}
        CPUは命令を処理し、ファイルや通信などの保護された機能をOSへ依頼します。

\par
\endgroup
\par
\endgroup


\section{通信経路}
\label{sec:07-practice-03}

実行中のブラウザは入力をHTTP要求へ組み、OSへ送信を依頼しました。名前解決、接続、暗号化、配送を経てサーバーへ着き、今度は逆順にアプリケーションまで上がります。各境界でデータへ何が付け加わり、何が取り外されるかを追います。

\begingroup
  \centering
  \begin{tikzpicture}[node distance=3mm,every node/.append style={font=\scriptsize}]
    \node[accentbox,minimum width=23mm] (browser) {クライアント};
    \node[box,minimum width=23mm,right=of browser] (dns) {DNS};
    \node[box,minimum width=23mm,right=of dns] (transport) {TCP／QUIC};
    \node[box,minimum width=23mm,right=of transport] (tls) {TLS};
    \node[accentbox,minimum width=23mm,right=of tls] (server) {サーバー};
    \draw[flow] (browser)--node[above]{名前}(dns);
    \draw[flow] (dns)--node[above]{IP}(transport);
    \draw[flow] (transport)--node[above]{接続}(tls);
    \draw[flow] (tls)--node[above]{HTTP}(server);
    \draw[softflow] (server.south) to[bend left=18] node[below]{応答は逆向きに戻る} (browser.south);
  \end{tikzpicture}
\begingroup
\paragraph{サーバーへ届くまで}
        DNSで宛先を求め、トランスポート接続を作り、TLSで相手と通信路を検証してからHTTP要求を送ります。再利用やキャッシュで一部を省く場合があります。

\par
\endgroup
\medskip
\begingroup
\paragraph{アプリへ渡るまで}
        ネットワークインタフェースからカーネルへ届いたバイト列を、OSが待機中のプロセスへ渡します。アプリがIPアドレスやポートを指定し、送受信を行うためにOSから受け取る通信の窓口が\term{ソケット}です。

\par
\endgroup
  \takeaway{「サーバーへ届かない」を一つの現象にまとめず、名前、経路、接続、TLS、HTTP、プロセスへ分解します。}
\par
\endgroup


\section{アプリケーション処理}
\label{sec:07-practice-04}

HTTPの本文が届くと、ただのバイト列は文字列、入力型、注文番号へ段階的に変換されます。ユースケースは注文と配送の事実を集め、モデルへ渡す文脈を作ります。ここでは設計図の箱が、実行時の順序として動き始めます。

\begingroup
  \centering
  \begin{tikzpicture}[node distance=2.4mm,every node/.append style={font=\scriptsize}]
    \node[accentbox] (bytes) {HTTPのバイト列};
    \node[box,right=of bytes] (parse) {解析\\構造化};
    \node[box,right=of parse] (validate) {検証\\境界で拒否};
    \node[box,right=of validate] (usecase) {ユースケース\\業務規則};
    \node[box,right=of usecase] (output) {応答へ変換};
    \draw[flow] (bytes)--(parse); \draw[flow] (parse)--(validate);
    \draw[flow] (validate)--(usecase); \draw[flow] (usecase)--(output);
  \end{tikzpicture}
  \vspace{5mm}
\begingroup
\paragraph{表現}
        JSONのバイト列を解析し、プログラム内の型、配列、スライス、マップ、オブジェクトへ変換します。外部表現を内部の値へ変えることを\term{デシリアライズ}と呼びます。内部の値を外部へ送れる形にすることを\term{シリアライズ}と呼びます。表現ごとに必要な容量と探索コストが異なります。

\par
\endgroup
\medskip
\begingroup
\paragraph{処理}
        分岐、探索、決めた順序へ要素を並べるソート、複数の値を合計など一つの結果へまとめる集約を組み合わせます。入力サイズと計算量から増え方を予測します。

\par
\endgroup
\medskip
\begingroup
\paragraph{設計}
        通信形式、業務規則、外部依存を別の責務に分け、変更理由が同じ処理をまとめます。

\par
\endgroup
\par
\endgroup


\section{メモリの使用}
\label{sec:07-practice-05}

この処理の各段階では、入力バッファ、解析後の値、注文データ、会話履歴、モデルへの要求が同時に存在します。どの値が一時的で、どれが要求を越えて共有され、いつ解放できるかを置き場所と寿命から見ます。

\begingroup
\begingroup
\paragraph{スタックとヒープ}
        \textbf{スタック}は関数呼び出しに対応する局所状態を保持します。
        \textbf{ヒープ}は呼び出しを越えて共有するオブジェクトを保持し、参照が残れば回収できません。

\paragraph{キャッシュとバッファ}
        後で再利用する結果を保つ領域が\term{キャッシュ}です。処理速度の異なる二者の間でデータを一時的に保つ領域が\term{バッファ}です。再計算や入出力回数を減らせますが、件数やバイト数の上限、いつ古いとみなすか、いつ解放するかの方針がなければ増え続けます。

\par
\endgroup
\medskip
\begingroup
\paragraph{使用量の見積もり}
        同時実行数$C$と一件当たりの平均使用量$m$から、合計使用量の目安を求めます。
        \centering
        $
          M_{\mathrm{total}} \approx
          M_{\mathrm{base}} + N_{\mathrm{concurrent}}M_{\mathrm{request}}
          + M_{\mathrm{cache}}
        $\par
        要素本体だけでなく、参照、型情報、容量の余り、一時コピーも含めます。

\paragraph{観測する量}
        プロセスが現在物理メモリ上に持つ量を\term{常駐メモリ}といいます。常駐メモリ、ヒープを型や用途別に分けた内訳、一秒当たりなどのメモリ確保量、GC時間、同時処理数、キャッシュ件数を、時間順に並べた\term{時系列}で比べます。

\par
\endgroup
\par
\endgroup


\section{機械学習モデルの組み込み}
\label{sec:07-practice-06}

注文と配送の事実を得た後、アプリはモデルへ文章化を依頼します。モデルは通常の関数のように呼べても、出力が確率的で、遅延と失敗の幅が大きく、事実性を自分では保証しません。設計上の境界と運用上の制御をここで重ねます。

\begingroup
  \centering
  \begin{tikzpicture}[node distance=3mm,every node/.append style={font=\scriptsize}]
    \node[accentbox] (case) {業務処理};
    \node[box,right=of case] (input) {入力整形\\機密情報を制御};
    \node[box,right=of input] (model) {モデル推論\\確率的な出力};
    \node[box,right=of model] (check) {形式・根拠・\\禁止事項を検査};
    \node[accentbox,right=of check] (decision) {採用／再試行／\\人へ引き継ぐ};
    \draw[flow] (case)--(input); \draw[flow] (input)--(model);
    \draw[flow] (model)--(check); \draw[flow] (check)--(decision);
  \end{tikzpicture}
  \vspace{6mm}
\begingroup
\paragraph{通常の関数と異なる点}
        同じ目的でも出力が変わり、流暢さと正しさは一致しません。タイムアウト、費用、モデル更新、分布変化も依存条件です。

\par
\endgroup
\medskip
\begingroup
\paragraph{アプリ側の責任}
        入力・出力の契約、評価集合、失敗時の代替、権限、監査、人の確認点を、モデルの外側で実装します。

\par
\endgroup
\par
\endgroup

\section{描画までの経路}

HTTP応答が成功しても、画面が更新されるとは限りません。
ブラウザはfetchの完了、JavaScriptの継続処理、DOM更新、スタイル計算、レイアウト、描画を順に進めます。
正常系では、Networkに成功応答が残り、Elementsに期待したDOMが現れ、Performanceに短い描画処理が記録されます。
異常系では、通信成功後の例外、更新対象の誤り、長いタスク、過剰な再レイアウトなどを分けて調べます。


\section{観測と仮説}
\label{sec:07-practice-07}

正常な一回を説明できると、異常な一回との差を探せます。外からの観測によって、システム内部で何が起きたかを調べられる性質を\term{可観測性}といいます。画面に見える「遅い」「答えが違う」を出発点にし、各層で測れる事実へ分け、その事実なら起こり得る原因を仮説として置きます。

出来事を時刻と説明付きで残した記録が\term{ログ}です。一つの要求が複数の処理やサービスを通る区間を時系列でつないだ記録が\term{トレース}です。件数、時間、使用量などを継続して数値化したものが\term{メトリクス}です。三者は答える問いが異なり、組み合わせて使います。

\begingroup
  \noindent
  \begin{tabularx}{\textwidth}{@{}>{\bfseries}p{19mm}YY@{}}
    \toprule
    階層 & 代表的な観測 & 読み取れること \\
    \midrule
    HTTP & ステータス、ヘッダー、本文、要求時間 & 要求の意味、応答の成否、待ち時間の位置 \\
    通信 & DNS時間、接続時間、再送、往復時間、TLS失敗 & 名前解決、経路、混雑、証明書、版の不一致 \\
    アプリ & ログ、トレース、処理件数、キュー長 & どの処理・依存先・入力で遅れたか \\
    OS & CPU、メモリ、GC、スレッド、入出力待ち & 計算、割り当て、競合、外部待ちのどれか \\
    コード & プロファイル、計算量、割り当て箇所 & 時間と空間を消費する関数・データ構造 \\
    AI & 入力群別の指標、根拠、費用、再現条件 & 確率的な失敗、偏り、分布変化、運用上の影響 \\
    \bottomrule
  \end{tabularx}
  \vspace{1mm}
  \takeaway{単一のログだけで結論を出さず、同じ時刻と入力条件、および一要求を追跡するために付けた要求IDの観測をつなぎます。通信の往復時間はRound-Trip Time（RTT）、入出力はInput/Output（I/O）とも表記します。}
\par
\endgroup


\section{遅延の調査}
\label{sec:07-practice-08}

五秒待ったという体験は、五秒かかる一つの処理があったことを意味しません。CPU計算、メモリ確保、キュー、DNS、接続、注文照会、モデル推論の時間を端から端までの流れに置き、どこで時間を失ったかを絞ります。

\begingroup
\begingroup
\paragraph{端から端までの時間}
        利用者が操作してから最終結果を受け取るまでの全体を\term{端から端まで}、英語でend-to-endといいます。応答時間は、待ち行列、通信、計算、保存装置や外部サービスからの返答待ちへ分解します。
        \[
          T_{\mathrm{total}}=
          T_{\mathrm{queue}}+T_{\mathrm{network}}+
          T_{\mathrm{compute}}+T_{\mathrm{I/O}}
        \]
        値を小さい順に並べた中央が\term{中央値}です。全体の$p$\%がその値以下になる境目を$p$\term{パーセンタイル}といいます。平均だけでは一部の非常に遅い要求が隠れるため、50、95、99パーセンタイル、入力サイズ、同時実行数ごとに分けます。

\paragraph{アルゴリズム}
        入力を2倍にしたときの変化を測ります。二重ループ、不要な全件走査、連結配列の再確保を疑います。

\par
\endgroup
\medskip
\begingroup
\paragraph{メモリとCPU}
        キャッシュミス、過剰なメモリ確保、コピー、GC、競合はCPU時間を増やします。実行時間やメモリ確保を関数・コード位置ごとに集計した記録を\term{プロファイル}といい、消費の大きい箇所を確認するために使います。

\paragraph{ネットワークと外部待ち}
        DNS、接続、TLS、最初の応答バイト、転送を別々に測ります。再試行は遅延と負荷をさらに増やし得ます。

\par
\endgroup
\par
\endgroup


\section{メモリ問題の調査}
\label{sec:07-practice-09}

メモリ使用量が増えたときも、すぐに「リーク」と名付けると道を誤ります。不要になったメモリをプログラムが参照し続けるなどして再利用できず、使用量が意図せず増え続ける問題が\term{メモリリーク}です。一件の値が大きいのか、同時に多く存在するのか、不要になっても参照が残るのかを、要求中のデータの寿命へ戻して分けます。

\begingroup
\begingroup
\paragraph{大きい}
        一件の配列、画像、応答、モデル入力が大きい場合です。全体を一度にメモリへ置かず、小さな部分を受け取りながら順に処理する\term{ストリーミング}、入力上限、より小さい表現を検討します。

\par
\endgroup
\medskip
\begingroup
\paragraph{多い}
        同時要求、オブジェクト数、一時コピーが多い場合です。並行数、バッファ再利用、データ構造を見直します。

\par
\endgroup
\medskip
\begingroup
\paragraph{長く残る}
        キャッシュ、グローバル参照、コールバック、元配列の共有で回収できない場合です。参照経路を調べます。

\par
\endgroup
  \vspace{5mm}
\paragraph{切り分けの順序}
    ある時点のヒープ上の値と参照を記録した\term{ヒープスナップショット}を、通常時と高負荷時で比較し、増えた型と確保元を特定します。その後、どの部品が値を管理するか、生存期間、上限、解放条件をコード上で確認します。

  \takeaway{GCがあることは、不要になった時点で自動的に解放されることを意味しません。到達可能な参照が残れば回収されません。}
\par
\endgroup


\section{通信障害の調査}
\label{sec:07-practice-10}

要求が届かない場合は、アプリの入口より前にも多くの失敗点があります。名前、経路、接続、証明書、HTTPのどこまで進んだかを観測し、最後に成功した境界と最初に失敗した境界を挟みます。

図の\term{NXDOMAIN}はDNSが「その名前は存在しない」と返す応答です。接続拒否は宛先へ到達したものの、そのポートで待つ処理がない場合などに返ります。タイムアウトは返答を得ないまま待ち時間の上限へ達した状態です。同じ「つながらない」でも、最後に得た応答によって到達地点が異なります。

\begingroup
  \centering
  \begin{tikzpicture}[node distance=2.6mm,every node/.append style={font=\scriptsize}]
    \node[accentbox] (dns) {名前をIPへ解決できたか};
    \node[box,right=of dns] (connect) {ポートへ接続できたか};
    \node[box,right=of connect] (tls) {TLS検証が通ったか};
    \node[box,right=of tls] (http) {HTTP応答を得たか};
    \draw[flow] (dns)--node[above]{はい}(connect);
    \draw[flow] (connect)--node[above]{はい}(tls);
    \draw[flow] (tls)--node[above]{はい}(http);
    \node[below=5mm of dns,align=center] (df) {NXDOMAIN、タイムアウト、\\古いキャッシュ};
    \node[below=5mm of connect,align=center] (cf) {拒否、経路なし、\\再送、ポート違い};
    \node[below=5mm of tls,align=center] (tf) {期限、名前、信頼、\\版、時計};
    \node[below=5mm of http,align=center] (hf) {4xx／5xx、形式、\\リダイレクト、容量};
    \draw[softflow] (dns)--node[left]{いいえ}(df);
    \draw[softflow] (connect)--node[left]{いいえ}(cf);
    \draw[softflow] (tls)--node[left]{いいえ}(tf);
    \draw[softflow] (http)--node[left]{いいえ}(hf);
  \end{tikzpicture}
  \vspace{5mm}
\paragraph{確認条件}
    発生時刻、接続元、対象名、IP、ポート、プロトコル版、証明書、HTTP要求を記録します。別の端末で成功しても、DNSキャッシュや経路が同じとは限りません。

\par
\endgroup

\section{描画の調査}

\term{Performance panel}は、ブラウザのJavaScript、スタイル、レイアウト、描画を時間軸で確認する画面です。
\term{flame chart}は、処理の入れ子と所要時間を帯で示す表示です。
\term{waterfall}は、ネットワーク要求の開始、待機、転送を横方向に並べた表示です。
\term{LCP}はLargest Contentful Paintの略で、主要コンテンツが描画された時点を示します。
\term{CLS}はCumulative Layout Shiftの略で、予期しない配置移動の累積を示します。

まずNetworkで応答の有無を確認します。
次にConsoleで例外を確認し、ElementsでDOM更新を確認します。
DOMが正しければ、Performance panelで長いタスク、再レイアウト、再描画を調べます。
各段階の観測が前段を反証できる順に進めます。


\section{変更容易性の改善}
\label{sec:07-practice-11}

障害を直せても、一つの変更が注文、通信、モデルの全体へ広がるなら、次の問題を生みます。修正時に触れた場所を手掛かりに、責任の混在、内部表現の漏れ、曖昧な失敗契約を探し、少しずつ境界を戻します。

\begingroup
\begingroup
\paragraph{兆候}
        一つの仕様変更で多数の場所を直す、テストに外部サービスが必須、条件分岐が各所に重複する、データ表現が境界を越えて漏れます。

\paragraph{構造上の原因}
        責務の混在、共有状態、暗黙の契約、循環依存、実装詳細への直接依存を、モジュール図と変更履歴で確認します。

\par
\endgroup
\medskip
\begingroup
\paragraph{小さな改善}
        まず現在の振る舞いをテストで固定し、境界を一つ選びます。インタフェースを狭め、必要な外部部品を内側で直接作らず外から渡す\term{依存性注入}によって、実装を交換できるようにします。一度に大きく動かさず、移動ごとに検証します。

\paragraph{効果測定}
        変更箇所数、テスト時間、障害範囲、理解に必要なモジュール数が減ったかを確認します。繰り返し現れる設計問題への典型的な解き方を\term{設計パターン}と呼びますが、名前の付いたパターンを導入すること自体を目的にしません。

\par
\endgroup
\par
\endgroup


\section{AI品質の調査}
\label{sec:07-practice-12}

文章が返ったことと、質問へ正しく答えたことは別です。どの注文データとプロンプトから、どのモデル版と生成条件で、どんな文が出たかを固定します。事実取得の誤りと文章生成の誤りを分けて初めて、改善先が見えます。

\begingroup
\begingroup
\paragraph{条件の固定}
        モデル版、プロンプト、アプリ提供者がモデルの役割や禁止事項を伝えるシステム指示、参照資料、温度、最大長、ツール結果、実行日時を残します。

\paragraph{失敗の分類}
        指示不遵守、形式違反、根拠不足、事実誤り、検索漏れ、拒否、遅延、費用超過を別の指標にします。

\par
\endgroup
\medskip
\begingroup
\paragraph{入力群別の測定}
        通常例だけでなく、境界例、曖昧な入力、長文、攻撃的入力、重要な利用者群を固定評価集合へ含めます。

\paragraph{利用境界}
        出力検査、根拠提示、再試行、代替処理、人の確認を加えます。改善が不十分なら自動化範囲を狭めます。

\par
\endgroup
  \takeaway{一つの成功例や失敗例ではなく、版管理した入力集合に対する分布として品質を比較します。}
\par
\endgroup


\section{障害調査の進め方}
\label{sec:07-practice-13}

個別の現象を見た後で、調査そのものの流れをまとめます。事実、仮説、観測、変更、再確認を混ぜずに順番へ置けば、思いつきを原因として扱わず、第三者が同じ道を再現できます。

\begingroup
  \centering
  \resizebox{\textwidth}{!}{%
  \begin{tikzpicture}[node distance=3mm,every node/.append style={font=\scriptsize}]
    \node[accentbox] (fact) {1. 現象を事実で記述};
    \node[box,right=of fact] (scope) {2. 範囲と時刻を限定};
    \node[box,right=of scope] (hypo) {3. 階層別に仮説化};
    \node[box,right=of hypo] (test) {4. 一変数ずつ検証};
    \node[accentbox,right=of test] (record) {5. 結果と次手を記録};
    \draw[flow] (fact)--(scope); \draw[flow] (scope)--(hypo);
    \draw[flow] (hypo)--(test); \draw[flow] (test)--(record);
    \draw[softflow] (record.south) to[bend left=18] node[below]{未解決なら更新} (hypo.south);
  \end{tikzpicture}%
  }
  \vspace{6mm}
\begingroup
\paragraph{事実}
        観測によって確かめた情報が\term{事実}です。「遅い」という評価だけでなく、対象、時刻、入力、期待値、実際に測った値、発生頻度、影響を書くと比較できます。

\par
\endgroup
\medskip
\begingroup
\paragraph{仮説}
        事実を説明する、まだ確定していない原因候補が\term{仮説}です。正しければ何が観測され、誤りなら何が観測されないかを先に決めます。どちらの結果でも仮説への判断が進む検査を選びます。

\par
\endgroup
\medskip
\begingroup
\paragraph{変更}
        影響を直ちに小さくする一時的な対応を\term{緊急回避}といいます。原因を取り除く対応を\term{恒久対策}といいます。二つを分けます。対策後も変更前と同じ指標を測り、別の問題を生んでいないか確認します。

\par
\endgroup
\par
\endgroup


\section{次の学習}
\label{sec:07-practice-14}

一つのリクエストを最後までたどると、自分が説明できない境界も見えてきます。その境界は欠点ではなく、次に深く学ぶ場所です。担当システムで実際に観測できる問いへ変え、必要な分野へ進みます。

表には、ここから先の道を示す用語も含まれます。OSが複数の処理へCPU時間を割り当てることが\term{スケジューリング}です。複数の処理が互いの解放を待ち続けて進めない状態が\term{デッドロック}です。アプリと必要なファイルを隔離した実行単位へまとめる仕組みを\term{コンテナ}と呼びます。

小さな部分問題の答えを保存し再利用して全体を解く方法が\term{動的計画法}です。データを速く探せるよう検索項目から位置への対応を別に持つ仕組みが\term{索引}です。複数のコンピュータが通信して一つの仕事を行う仕組みを\term{分散システム}といいます。相手が誰かを確かめる認証に対し、その相手が何をしてよいか決める処理が\term{認可}です。想定する攻撃者、守る対象、攻撃経路、対策を整理することを\term{脅威分析}といいます。

データの集め方、学習、評価、配備、監視、更新を継続して管理する実務を\term{MLOps}と呼びます。Machine Learningと、開発・運用を協力させるOperationsを組み合わせた語です。

\begingroup
  \noindent
  \begin{tabularx}{\textwidth}{@{}>{\bfseries}p{30mm}YY@{}}
    \toprule
    分野 & 次の主題 & 実務で答えられる問い \\
    \midrule
    アーキテクチャ & パイプライン、キャッシュ、並列性、性能測定 & 同じコードでもCPUや配置で性能が変わるのはなぜか \\
    OS・並行処理 & スケジューリング、同期、I/O、ファイル、コンテナ & 待ち、競合、デッドロック、資源上限をどう調べるか \\
    アルゴリズム & グラフ、動的計画法、索引、確率的手法 & 入力規模に対して処理時間とメモリをどう抑えるか \\
    設計・テスト & アーキテクチャ、契約、テスト戦略、継続的改善 & 変更を安全に小さく届ける境界をどう作るか \\
    通信・安全性 & ルーティング、分散、認証・認可、脅威分析 & 信頼境界を越えるデータと権限をどう守るか \\
    機械学習・AI & 統計、最適化、評価、MLOps、安全な運用 & 予測品質をどう測り、変化と失敗をどう管理するか \\
    \bottomrule
  \end{tabularx}
  \vspace{1mm}
  \takeaway{用語の暗記ではなく、自分の担当システムで再現できる小さな観測・実装・説明を学習単位にします。}
\par
\endgroup


\section{もう一度たどる：注文確認AI}
\label{sec:07-practice-15}

最後に、冒頭と同じ要求を自分の手で組み立て直します。今度は完成した説明を読むのではなく、各層の出力が次の入力になるように線をつなぎ、正常時と異常時の両方を一枚の物語へします。

\begingroup
\begingroup
\paragraph{シナリオ}
        利用者がブラウザへ「注文番号1234はいつ届きますか」と入力します。Web APIは注文を検索し、配送サービスから到着予定を取得します。その事実を文脈へ含めて生成AIを呼び出し、利用者向けの応答を返します。

\paragraph{発生している現象}
        通常は800ミリ秒ですが、遅い側の1\%、すなわち99パーセンタイルが5秒を超えます。ミリ秒は千分の一秒です。多数の要求を送り、想定する利用量で性能と安定性を確かめる\term{負荷試験}では、メモリが増え続けます。まれにTLSエラーも発生し、応答には配送サービスが返していない日時が混じることがあります。

\par
\endgroup
\medskip
\begingroup
\paragraph{前提としてよいこと}
        入力はHTTPSのJSONです。アプリはコンパイル済みプロセスとして動作し、会話履歴をスライス、注文索引をマップで保持します。

\paragraph{演習の目的}
        すぐに原因を断定するのではなく、各階層の処理、失敗点、必要な観測、検証順序を説明します。

\par
\endgroup
\par
\endgroup


\section{正常な旅}
\label{sec:07-practice-16}

まず、すべてが成功する一回を途中で飛ばさずに描きます。矢印には単に「呼ぶ」と書かず、渡す表現と受け取る側の保証を添えてください。

\begingroup
  \begin{enumerate}
    \setlength{\itemsep}{1mm}
    \item ソースコードからプロセス起動までを、コンパイラ、機械語、OS、仮想メモリ、CPUを使って説明してください。
    \item ブラウザからAPIへ要求が届くまでを、DNS、IP、TCPまたはQUIC、TLS、HTTPの順で図示してください。
    \item JSONが型付きの質問へ変換され、マップから注文を探し、配送予定と会話履歴をモデルの文脈へ組み、JSON応答になる流れを書いてください。
    \item 処理中の主要データを、スタック、ヒープ、共有キャッシュのどこへ置くか仮定し、生存期間を添えてください。
    \item AIによる文章化を配送事実の取得から分離し、入力制御、タイムアウト、出力検査、失敗時の代替を境界へ書いてください。
  \end{enumerate}
  \vspace{1mm}
\paragraph{提出物}
    A3用紙一枚相当のシーケンス図または階層図と、各矢印を説明する800〜1,200字の文章を作ります。図中の矢印には、渡すデータと提供する保証を書きます。

\par
\endgroup


\section{途切れた旅}
\label{sec:07-practice-17}

次に、同じ経路のどこかが遅れ、あふれ、届かず、誤る場合を扱います。現象ごとに旅全体を作り直すのではなく、正常系との差が最初に現れる境界を探します。

\begingroup
  \noindent
  \begin{tabularx}{\textwidth}{@{}>{\bfseries}p{22mm}YY@{}}
    \toprule
    現象 & 少なくとも立てる仮説 & 必ず含める観測 \\
    \midrule
    遅延 & 計算量、割り当て、外部待ち、キューの各観点 & 入力サイズ別時間、プロファイル、依存先時間、上位百分位 \\
    メモリ増加 & 一件が大きい、多数存在する、参照が残るの各観点 & 型別ヒープ差分、確保元、同時数、GC、キャッシュ件数 \\
    TLS失敗 & 名前、経路、時刻、証明書、プロトコル版の各観点 & 発生元、IP、証明書鎖、有効期限、時計、エラー全文 \\
    AI誤り & 入力、検索、モデル出力、検査、運用境界の各観点 & 固定評価集合、根拠一致率、条件、モデル版、失敗時の扱い \\
    \bottomrule
  \end{tabularx}
  \vspace{4mm}
\paragraph{計画の書式}
    各現象について「事実 → 仮説 → 仮説を支持する観測 → 反証する観測 → 最初の安全な検証 → 対策後の確認指標」を一組にします。

\par
\endgroup


\section{評価方法}
\label{sec:07-practice-18}

説明の良さは、用語の数ではなく、第三者が同じ経路をたどれるかで評価します。層を網羅し、矢印の因果を示し、事実と仮定を分け、調査で反証できることを基準にします。

\begingroup
  \noindent
  \begin{tabularx}{\textwidth}{@{}>{\bfseries}p{25mm}p{14mm}Y@{}}
    \toprule
    評価観点 & 配点 & 達成基準 \\
    \midrule
    階層の網羅 & 20 & 言語、OS、メモリ、CPU、通信、アプリ、AIの役割を区別しています。 \\
    因果関係 & 25 & データと制御の流れを矢印で示し、前段の出力が次段でどう使われるか説明しています。 \\
    技術的正確さ & 20 & 各層が保証することと保証しないこと、平均と最悪、学習と推論を混同していません。 \\
    調査計画 & 25 & 仮説を観測で支持・反証でき、条件、比較対象、安全な検証順序が明確です。 \\
    伝達品質 & 10 & 用語を定義し、事実・仮定・推測を分け、第三者が同じ調査を再現できます。 \\
    \bottomrule
  \end{tabularx}
  \vspace{1mm}
\paragraph{相互レビュー}
    ペアで説明を交換し、矢印を一つ選んで「誰が、何を、どの表現で、どの保証の下に渡すか」を質問します。答えられない境界を追加調査の候補にします。

\par
\endgroup


\section{送信ボタンへ戻る}
\label{sec:07-practice-19}

もう一度、利用者の画面を見ます。「注文番号1234はいつ届きますか」という一文は、入力装置の電気信号から始まりました。CPUが命令を実行し、OS上のプロセスが文字列を値として保持し、設計された境界に沿って注文と配送の事実を集めました。HTTP要求はTLSで守られてネットワークを渡り、モデルは渡された事実を含む文脈から応答を生成しました。その文は同じ層を戻り、ブラウザが画面へ描きました。

どの層も、一つ前の層の仕事をなかったことにはしていません。上位は下位の細部を隠しながら、その計算、記憶、配送を使っています。だから問題が起きたときは、抽象化を捨ててすべてを調べるのではなく、最後に成功した境界まで戻り、そこから一段ずつ降ります。

この物語に終点はありますが、システムには次の要求が届きます。新しい機能も障害も、同じようにデータの姿、受け渡す相手、境界の保証をたどれば、どこから理解を始めるべきかが見えます。

\takeaway{物理からアプリまでをつなぐ鍵は、全詳細を同時に覚えることではありません。各層でデータがどんな姿を取り、何を保証されて次へ渡るかを、途切れない因果として説明することです。}

\subsection*{第9章の索引用語}

\begin{description}
  \item[\term{全体像}] 一要求に関係する層、境界、観測点を一枚にまとめた地図です。
  \item[\term{観測点}] 値、状態、時間、失敗を記録できる境界です。
  \item[\term{span}] 一つの処理区間と属性、開始、終了を表すトレースの単位です。
  \item[\term{trace\_id}] 複数のspanを一要求のトレースへ結び付ける識別子です。
  \item[\term{correlation id}] 異なるログやメッセージを同じ処理へ関連づける識別子です。
  \item[\term{OpenTelemetry}] トレース、メトリクス、ログの計装と伝送を標準化するプロジェクトです。
  \item[\term{SLI}] Service Level Indicatorの略で、サービス水準を測る指標です。
  \item[\term{SLO}] Service Level Objectiveの略で、SLIに対して定める目標です。
  \item[\term{エラーバジェット}] 目標期間内でSLOを満たしつつ許容できる失敗量です。
  \item[\term{プロファイリング}] CPU時間やメモリ割り当ての分布を測る分析です。
  \item[\term{ボトルネック}] 全体の処理量や時間を主に制約する箇所です。
  \item[\term{検証}] 仮説や修正結果を観測によって確かめる行為です。
  \item[\term{再現}] 同じ条件と手順で同じ症状を再び観測することです。
  \item[\term{切り分け}] 観測と反証によって原因候補の範囲を狭めることです。
  \item[\term{DNS委任}] 名前空間の一部を別の権威DNSへ任せる関係です。
  \item[\term{経路障害}] パケットを宛先へ届ける経路が利用できない、または不安定な状態です。
  \item[\term{TLS}] Transport Layer Securityの略で、通信の機密性、完全性、認証を支えます。
  \item[\term{AI品質}] 利用目的に対するモデル出力の正確さ、安全性、安定性などの性質です。
  \item[\term{評価データ}] 品質指標を測るため、学習とは分けて用意する入力と期待結果です。
  \item[\term{インシデント}] サービスへ望ましくない影響が生じ、対応と記録が必要な出来事です。
  \item[\term{ポストモーテム}] インシデント後に事実、影響、原因、改善を非難ではなく学習目的で記録する振り返りです。
  \item[\term{正常系}] 想定した前提と入力で期待する処理が完了する流れです。
  \item[\term{異常系}] 失敗、例外、制限超過など、通常と異なる処理へ進む流れです。
\end{description}

